This paper presents an evidence-grounded English manuscript draft for a simulation-oriented planning pipeline built around three implemented components: Memento-style case retrieval, a Hope adaptive weighting controller, and a reflection loop for iterative prompt refinement. The system generates structured joint-operation plans, evaluates them through a five-stage kill-chain simulation model, and exports the selected plan into DoDAF OV-5b, DoDAF OV-6c, and C2SIM artifacts. To avoid unsupported claims, the draft only describes mechanisms that are explicitly present in the current codebase and experiment outputs and frames the study as an application-oriented integration paper rather than as a claim of a new standalone planning theory.

The main single-scenario ablation study uses the repository's sample coastal joint-assault scenario with 10 optimization iterations, 50 Monte Carlo simulation runs, a fixed random seed of 7, and writeback disabled. Under optimized HOPE parameters ($\theta=0.9$, $\epsilon=0.005$), the full pipeline improves mission success from 0.6260 to 0.6707, overall effectiveness from 0.7328 to 0.8011, and the loss-exchange ratio from 1.6564 to 2.2721 relative to the pure large language model baseline. A multi-seed replication across five seeds confirms the Full pipeline's gain is statistically significant under default parameters (paired $t(4)=5.593$, $p<0.01$, Cohen's $d=2.50$). A supplementary multi-seed run with optimized parameters ($\theta=0.9$, $\epsilon=0.005$, $n=3$ seeds) yields a mean MS gain of +0.0606 over Pure LLM, confirming the optimized gain directionally generalizes across seeds, though seed-specific variation suggests that the optimal parameter set may itself be seed-dependent. The main cross-scenario transfer study uses five scenario families and leave-one-source-scenario-out case banks derived from 250 source cases. In that evidence set, the full pipeline outperforms the pure baseline in all five scenarios, increasing average mission success from 0.6595 to 0.6828 and average overall effectiveness from 0.7417 to 0.7924, while the relative ordering among the Memento-only and Hope-only variants remains scenario-dependent.

The draft also documents the remaining validity limits, including single-seed evidence, scenario dependence in module gains, and the lack of external baselines or significance testing. In particular, the reported confidence intervals reflect within-seed Monte Carlo variance under the fixed seed rather than across-seed variance. The refreshed canonical result bundles now serialize a top-level runtime manifest that records the local backend identifier, software stack, available CPU/GPU metadata, and wall-clock timing summaries for the reported evidence.
